\chapter{Analytical proposition and evaluation questions}
\label{cap:hipotesis}

The study distinguishes an algebraic consequence of the swap rule from two questions that require simulation. These are claims about the specified zero-fee micro-world; they are not hypotheses about observed markets.

\section{H1: ideal correction under bounded arbitrage}

Let $e_t=\log(p_t^D/p_t^C)$ and $D_t=(|e_t|-\gamma)_+$. In an isolated, zero-fee pool with a fixed reference, no other orders and no binding inventory limit, the maximum order $q=cx$ changes the log price by $\delta_c=2\log(1+c)$. Since the arbitrageur stops at the band, the out-of-band error evolves as
\[
  D_{t+1}=\max(0,D_t-\delta_c),\qquad
  N_\gamma=\left\lceil\frac{D_0}{2\log(1+c)}\right\rceil \quad(c>0).
\]
For $c=0$, no finite correction occurs. This is a mechanical benchmark derived from the constant-product swap equation, not an empirical finding. Simulation is useful only for quantifying departures from this benchmark caused by other agents, event order, finite horizon and the specified inventory rules.

\section{H2: withdrawal, order scale and tracking response}

The available experiments address a narrower evaluation question:
% QUESTION_LOCK:h2:BEGIN
\emph{Within the tested capacity and depth levels, how does the specified one-time liquidity-withdrawal rule affect post-shock tracking, and how does that effect vary when order sizes scale with reserves or panic behavior is disabled?}
% QUESTION_LOCK:h2:END
The E3 rule is a one-time threshold trigger, not an estimated adaptive policy: it withdraws 45\,\% of reserves as soon as $|e_t|>0.02$. The executed shocks exceed this threshold in every E3 run. The original peak displacement from the pre-shock AMM price is treated as a diagnostic only; it is not a measure of overshoot relative to the new reference.

\section{H3: descriptive accuracy of reduced maps}

The map clause of the research question is operationalized as predictive description of continuous tracking-error paths on held-out simulator runs. The finite-window trajectory labels summarize those paths; they are not independently observed regimes or targets in a classification test. H3 asks:
% QUESTION_LOCK:h3:BEGIN
\emph{How accurately does the two-state map predict held-out tracking-error paths relative to persistence and a scalar AR(1), and how does a band-aware map perform against the affine map in recursive replay?}
% QUESTION_LOCK:h3:END
The comparison concerns simulator-generated paths. The hybrid replay receives the realized future CEX path, so its result is not a real-time forecast test. No regime-classification or stability hypothesis is advanced.
